\section{Математические основы RSA}

Безопасность алгоритма RSA базируется на вычислительной сложности задачи факторизации больших целых чисел. Генерация ключей начинается с выбора двух различных простых чисел \texttt{p} и \texttt{q}. Первым шагом вычисляется их произведение:
$n = p \cdot q$

Число \texttt{n} называется модулем и используется как часть как открытого, так и закрытого ключей. Далее вычисляется функция Эйлера от числа \texttt{n}:
\begin{equation}
\label{eq:euler}
\varphi(n) = (p - 1)(q - 1)
\end{equation}

После этого выбирается целое число \texttt{e}, такое что $1 < e < \varphi(n)$, причем \texttt{e} и $\varphi(n)$ должны быть взаимно простыми. Число \texttt{e} становится открытой экспонентой. Затем вычисляется число \texttt{d}, называемое секретной экспонентой, которое удовлетворяет следующему условию:
\begin{equation}
\label{eq:secret_d}
d \cdot e \equiv 1 \pmod{\varphi(n)}
\end{equation}

Процессы шифрования исходного сообщения \texttt{m} и расшифрования полученного шифротекста \texttt{c} можно описать с помощью следующих математических выражений:
\begin{align}
c &\equiv m^e \pmod{n} \label{eq:encrypt} \\
m &\equiv c^d \pmod{n} \label{eq:decrypt}
\end{align}

Как видно из уравнений~\ref{eq:encrypt} и~\ref{eq:decrypt}, а также свойства функции Эйлера~\ref{eq:euler}, процесс восстановления данных требует наличия секретного ключа.
